\section{Related Work}
\label{sec:related-work}

\subsection{Circuit simulators}

SPICE and its descendants~\cite{nagel1973spice} remain the reference implementation of
circuit simulation, built around modified nodal analysis~\cite{ho1975mna}: node voltages and
a small number of auxiliary branch currents (one per independent voltage source or inductor,
in the classical formulation) are solved for simultaneously at each timestep, with reactive
elements replaced by companion resistive/source models derived from a numerical integration
rule. Educational tools such as the Falstad circuit simulator and PhET's Circuit
Construction Kit lower the barrier to entry by replacing SPICE's netlist syntax with direct
schematic manipulation, and are widely used in introductory courses. All three, however,
present the circuit as a two-dimensional schematic. Mine Memristors uses the same underlying
numerical technique -- Sections~\ref{sec:architecture} and~\ref{sec:circuit-solver} describe
an MNA solver with trapezoidal-integration companion models for reactive elements, the same
family of technique SPICE uses -- but embeds it in a three-dimensional, first-person, block-based
world instead of a schematic canvas.

\subsection{Minecraft in education}
\label{sec:minecraft-in-education}

The instructional value of games and simulations is well studied in general: a meta-analysis
by Sitzmann~\cite{sitzmann2011meta} found simulation games to be more effective for training
self-efficacy and knowledge retention than comparison instructional methods across a range
of technical domains. Minecraft occupies an unusually large share of that broader
literature specifically, to the point that Microsoft maintains and sells a dedicated,
classroom-oriented variant, \emph{Minecraft: Education Edition}, aimed squarely at schools --
itself a useful indicator of how far the game's informal adoption by teachers had already
progressed before a vendor built a product around it. Nebel et al.'s
review~\cite{nebel2016mining} surveys this literature broadly and identifies Minecraft's
core appeal for education as the same thing that makes it appealing as a game: an open,
constructive world in which a learner's own artifact, not a worked example, is the object of
study, closely matching Papert's constructionist framing~\cite{papert1980mindstorms}
discussed further in Section~\ref{sec:pedagogy}.

Published case studies span a wide range of subjects well outside engineering. Bos et
al.~\cite{bos2014mathematics} describe using Minecraft to teach elementary-school
mathematics concepts such as area, perimeter, and ratio through building tasks; Short's
account~\cite{short2012teaching} covers using the game to teach general science concepts in
a secondary-school setting; and Cipollone et al.~\cite{cipollone2014creative} present a case
study of Minecraft used as an open-ended creative tool, with an emphasis on collaboration
and planning skills rather than any specific curricular content. Callaghan's broader
survey~\cite{callaghan2016investigating} of Minecraft's role across primary and secondary
classrooms echoes a theme common to all of the above: teachers report high engagement and
motivation fairly consistently, but rigorously measured gains in specific learning outcomes
are reported less often and less consistently, a caveat also present in the general
simulation-games literature~\cite{sitzmann2011meta} and one this paper does not claim to
have resolved for Mine Memristors specifically (Section~\ref{sec:limitations}).

Within that broader literature, the subset closest to this paper's subject is Minecraft's
native redstone subsystem, which is already a common entry point for teaching digital logic
and computer architecture concepts: players build AND/OR/NOT gates, adders, and working
memory cells entirely out of in-game blocks, and Dezuanni et
al.~\cite{dezuanni2015redstone} specifically document children describing and reasoning
about redstone circuits using explicitly electrical language, even though the underlying
mechanic is a discrete, boolean signal rather than a continuous one (discussed further, with
the actual mechanics of redstone, in Section~\ref{sec:minecraft-mechanics}). This is the
literature Mine Memristors sits closest to and departs from most directly: the case studies
above use Minecraft's \emph{existing} mechanics, redstone included, as the vehicle for a
lesson, whereas Mine Memristors adds an entirely new, continuous-time simulation layer
because redstone's discrete signal has no way to represent voltage, current, or reactive
state at all. To the author's knowledge, it is the first Minecraft modification to embed a
real continuous-time circuit solver for this purpose, rather than approximating one with the
game's existing discrete redstone mechanics.

\subsection{Memristors as a teaching subject}

The memristor was postulated on symmetry grounds by Chua~\cite{chua1971memristor} as a
fourth fundamental two-terminal circuit element, relating charge and flux, alongside the
resistor, capacitor, and inductor; a physical device matching the predicted
charge-controlled behavior was reported by Strukov et al.\ at
HP~Labs~\cite{strukov2008missing}. Because the memristor is both the newest of the four
canonical elements and the one whose behavior (state that persists and depends on the
history of current through the device) is least like anything in a first circuits course,
it is a natural stress-test for whether a teaching tool actually solves the underlying
mathematics rather than approximating it -- an in-game device whose resistance visibly
drifts with accumulated charge, and whose state survives being disconnected and reconnected,
is a much more concrete introduction to the concept than the differential equation alone.
